\section{Consequences for named inequality families}
\label{sec:classes}

The base construction transfers to several infinite families because its
rounded psd violations are exactly its hypermetric violations. A second,
orthogonal extension permits coefficients of absolute value at most one
without sacrificing positive definiteness.

\begin{theorem}\label{thm:binary-families}
Separation of each of the following families is strongly $\NP$-complete:
hypermetric correlation inequalities, gap-$1$ inequalities, rounded psd
inequalities, Boros--Hammer inequalities, and the union of the $k$-gonal
families over all $k$.
For every fixed $K\geq3$, hardness holds at rational distance vectors
$d\in\MET(V)$ with $Z(d)\in\operatorname{relint}\ELL(V)$ that satisfy
every rounded psd inequality of gonality at most $2K-3$.
Their Boolean quadric moment matrices are positive definite, and their
coordinates have polynomially bounded numerators and a polynomially
bounded common denominator.
\end{theorem}

\begin{proof}
Use the points $\varepsilon d$ of
\cref{thm:relaxation-hardness}, with $q\geq K$.
Every hypermetric inequality is a gap-$1$ inequality: its coefficient sum
is one, so its gap is at most one, while parity excludes zero.
Every gap-$1$ inequality is rounded psd. At these points, the only
rounded psd violations have coefficient sum $\pm1$ and correspond to
exact covers. The members of the $k$-gonal union with even coefficient
sum are psd inequalities and therefore hold.
Thus each family has a violation if and only if an exact cover exists.
The covariance map identifies hypermetric inequalities with hypermetric
correlation inequalities and rounded psd inequalities with Boros--Hammer
inequalities. The size bounds and relaxation properties follow from
\cref{thm:relaxation-hardness}.

For membership in $\NP$, hypermetric correlation inequalities use
\cref{lem:short-hypermetric-witness}. Rounded-psd and Boros--Hammer
inequalities use \cref{lem:short-split-witness} and their exact covariance
and split correspondences.
For the $k$-gonal union, an even-sum violation requires an indefinite
$Z=J-2D$. By \cref{lem:rational-negative-direction}, it has an integral
negative direction of polynomial encoding length; doubling that direction
gives an even-sum certificate. An odd-sum violation
is rounded psd.

For gap-$1$, suppose $b$ is violated and choose $a\in\{\pm1\}^{V}$ with
$\abs{a\trans b}=1$. After a possible global negation, the vector
$b'=\diag(a)b$ has coefficient sum one. Switch $d$ on
$S=\{u:a_u=-1\}$, so that $Z(d^S)=\diag(a)Z(d)\diag(a)$.
Then $(b')\trans Z(d^S)b'<1$, equivalently $Q(b',d^S)>0$.
Conversely, a hypermetric violation on a switched input gives a gap-$1$
violation on the original input: the switching signs witness a signed
sum of absolute value one, and odd parity proves that the gap is exactly
one. A certificate consists of $S$ and a short hypermetric witness for
$d^S$. Its length is polynomial by
\cref{lem:short-hypermetric-witness}. This argument applies to arbitrary
rational inputs, including inputs outside $[0,1]^{E(V)}$.
\end{proof}

\subsection{Signed Boolean quadric families and their facets}

Recall the signed family \eqref{eq:finite-BH}, whose slack is
$h_Y(\one_S-\one_T,s)/2$. Its cut subfamily has $s=0$; its clique
subfamily has $T=\varnothing$ and $0\leq s\leq\abs S-1$.
These are the families (10)--(12) in \citet{Letchford2022}; the
unrestricted Boros--Hammer family is (13). See also
\citet{Padberg1989} and \citet{LetchfordSorensen2012}.

Complement one variable $g$: set
\[
 x'_g=1-x_g,\qquad y'_{ig}=x_i-y_{ig}\quad(i\ne g),
\]
and leave the other coordinates unchanged. Denote this map by $\psi_g$.
Its moment matrix is $AYA\trans$, where $A$ is the identity except that
row $g$ is $e_0\trans-e_g\trans$. Since $A^2=I$, it preserves positive
definiteness. If $v'$ is $v$ with its $g$-entry negated, then
\begin{equation}\label{eq:bh-switching}
 h_{AYA\trans}(v,s)=h_Y(v',s-v_g).
\end{equation}
Indeed, with $\widehat u=(-s,v)$,
$h_Y(v,s)=\widehat u\trans Y\widehat u-e_0\trans Y\widehat u$,
$A\trans\widehat u=(-(s-v_g),v')$, and $e_0\trans A=e_0\trans$.
The shift of $s$, as well as the sign change of $v_g$, is part of the
switching rule. On the Boolean cube this map complements $x_g$, so it
is an affine automorphism of the Boolean quadric polytope. On its cut
image it is switching on $\{g\}$.

We prove the facet property needed below directly.

\begin{lemma}\label{lem:clique-facets}
Let $S\subseteq W$, $t=\abs S\geq3$, and $1\leq r\leq t-2$.
The clique specialization of \eqref{eq:finite-BH} with parameter $r$
defines a facet of $\BQP(W)$, including when $S\ne W$.
\end{lemma}

\begin{proof}
At a binary point its slack is
\[
 F(x)=\sum_{\{i,j\}\subseteq S}x_ix_j-r\sum_{i\in S}x_i+
       \frac{r(r+1)}2
     =\frac{(\one_S\trans x-r)(\one_S\trans x-r-1)}2.
\]
Its tight points are exactly those with $\one_S\trans x\in\{r,r+1\}$.
An affine equation on Boolean quadric coordinates is a multilinear
polynomial of degree at most two on the Boolean cube. First suppose
such a polynomial
$p=c+\sum_{i\in S}a_ix_i+\sum_{\{i,j\}\subseteq S}a_{ij}x_ix_j$
vanishes on both tight layers.
For distinct $i,j\in S$, comparing the points supported on
$R\cup\{i\}$ and $R\cup\{j\}$ gives
\[
 a_i-a_j+\sum_{k\in R}(a_{ik}-a_{jk})=0
\]
for every $R\subseteq S\setminus\{i,j\}$ of size $r-1$ or $r$.
For any $k\notin\{i,j\}$, choose $R$ of size $r-1$ excluding $k$;
this is possible because $r\leq t-2$.
Subtract the equations for $R$ and $R\cup\{k\}$. Thus
$a_{ik}=a_{jk}$. All pair coefficients have a common value $a$,
and all linear coefficients have a common value $b$.
The two layer equations give $b=-ra$ and $c=r(r+1)a/2$.
Therefore $p=aF$.

Now allow variables outside $S$. Set them all to zero. The remaining
polynomial is a multiple of $F$. Setting one outside variable to one
shows that its coefficient, an affine function on $S$, vanishes on
both layers. Such an affine function is zero: swaps show that all
its linear coefficients agree, and the two distinct layer values
then force that coefficient and the constant to vanish.
Setting two outside variables to one eliminates their pair
coefficient. Hence every affine equation on the tight points is a
multiple of $F$. The Boolean quadric polytope has full dimension:
evaluation at $0,e_i,e_i+e_j$ shows that $1,x_i,x_ix_j$ are linearly
independent on the Boolean cube.
Its tight points therefore span a face of codimension one.
\end{proof}

\begin{theorem}\label{thm:finite-families}
Separation of Padberg's cut inequalities, Padberg's clique inequalities,
and \eqref{eq:finite-BH} is strongly $\NP$-complete.
For every fixed $K\geq3$, hardness holds at positive definite Boolean
quadric moment points whose cut images $d$ satisfy
$d\in\MET(V)$, $Z(d)\in\operatorname{relint}\ELL(V)$, and every
rounded psd inequality of gonality at most $2K-3$.
Their coordinates have a polynomially bounded common denominator and
polynomially bounded numerators. Every violated inequality in these
constructions defines a facet.
\end{theorem}

\begin{proof}
At the point $Y_\varepsilon$ of \cref{thm:relaxation-hardness},
\eqref{eq:BH-rounded} and the exact list of split violations give precisely
the violated Boros--Hammer data
\[
 (v,s)=(\one_C-e_g,0),\qquad(-\one_C+e_g,-1),
\]
for the exact covers $C$. These are two forms of the same inequality.
Consequently, up to $(v,s)\sim(-v,-s-1)$, the violated members of
\eqref{eq:finite-BH} are exactly $S=C,T=\{g\},s=0$.
They are cut inequalities, with violation $1/(4N)$.

At $\psi_g(Y_\varepsilon)$, \eqref{eq:bh-switching} gives exactly
\[
 (v,s)=(\one_C+e_g,1),\qquad(-\one_C-e_g,-2).
\]
Thus the violated clique inequalities are exactly
$S=C\cup\{g\},s=1$, again with violation $1/(4N)$.
They are also the only violated members of \eqref{eq:finite-BH},
up to the same equivalence.
Their supports have size $q+1$, so \cref{lem:clique-facets} applies.
The cut violators are their images under $\psi_g$ and are therefore
facets too.

The switched moment matrix is positive definite by congruence. Its changed
entries are
\[
 x'_g=1-\frac{2p-1}{4N},\qquad y'_{ig}=\frac{11}{2N}
\]
for X3C, so its coordinates are in $[0,1]$ with common denominator $4N$.
Cut switching is diagonal-sign congruence on $Z$ and preserves the
gonality of rounded psd inequalities. It permutes triangle and perimeter
inequalities, hence preserves the metric polytope as well as the
elliptope interior. All promised properties follow from
\cref{thm:relaxation-hardness}. For a fixed $K$, the X3C restriction
$q\geq K$ remains strongly $\NP$-hard, so these encodings establish
the claimed strong hardness.

For membership in $\NP$, a cut inequality is specified by $(S,T)$.
A clique inequality is specified by $S$ and an integer in a range of
length $\abs S$. For \eqref{eq:finite-BH}, fix $(S,T)$ and put
$\ell=x(S)-x(T)$. Its violation as a function of $s$ is
$-s^2/2+s(\ell-1/2)+c$, for a constant $c$.
Its maximum over the integers is attained at $s=\lfloor\ell\rfloor$,
which has polynomial bit length. Thus $(S,T)$ alone suffices as a
certificate: the verifier computes this parameter and the slack.
\end{proof}

\subsection{Pure hypermetric and cut-form clique families}

The odd-clique family in \cref{tab:classes} is precisely the set of
switchings of the all-positive clique inequalities
\begin{equation}\label{eq:cut-clique}
 \sum_{\{u,v\}\subseteq S}d_{uv}\leq\frac{\abs S^2-1}{4},
 \qquad \abs S\text{ odd},
\end{equation}
usually called the Barahona--Mahjoub clique inequalities; see
\citet{BarahonaMahjoub1986,Letchford2022,GKL2012}.

Take the X3C construction with $q\geq3$, and put
\[
 m=q-2,\qquad \tau^2=\frac1{8q},\qquad
 \widetilde M=M\oplus\tau^2 I_m.
\]
Add indices $o_1,\ldots,o_m$ and let
$\widetilde V=V\cup\{o_1,\ldots,o_m\}$.
Its covariance distance $\widetilde d=d(\widetilde M)$ satisfies
\begin{align*}
 \widetilde d_{0o_j}&=\tau^2,&
 \widetilde d_{o_jo_k}&=2\tau^2\quad(j\ne k),&
 \widetilde d_{o_ju}&=d_{0u}+\tau^2\quad(u\in W).
\end{align*}
In a Gram representation, the new points are the root moved by $\tau$
in distinct new orthogonal directions. They are distinct from each other
and from the original points.
Let $\widetilde Y_\varepsilon$ be the scaled moment matrix of this
extension, with the original $N=8n+2p+1$ and $\varepsilon=1/N$.
Let $d''$ be $\varepsilon\widetilde d$ switched on
$U=\{g,o_1,\ldots,o_m\}$.

\begin{theorem}\label{thm:pure-clique-hardness}
Separation of pure hypermetric inequalities, odd-clique inequalities,
and the all-positive clique inequalities \eqref{eq:cut-clique} is strongly
$\NP$-complete.
For every fixed $K\geq3$, hardness holds at rational distance vectors
$d\in\MET(V)$ with $Z(d)\in\operatorname{relint}\ELL(V)$ that satisfy every rounded psd
inequality of gonality at most $2K-3$. Their moment and covariance
Gram matrices are positive definite, and their coordinates have
polynomially bounded numerators and common denominator.

More precisely, at $\varepsilon\widetilde d$ the pure hypermetric
violators for a cover $C$ are exactly
\begin{equation}\label{eq:pure-violators}
 b^*=(0,\one_C,-1,-\one_m),\qquad
 b^{(j)}=(-1,\one_C,-1,-\one_m+e_j),\quad 1\leq j\leq m,
\end{equation}
in the order $(0,1,\ldots,n,g,o_1,\ldots,o_m)$.
The odd-clique violators are these vectors and their negatives.
At $d''$, the violated all-positive clique inequalities are exactly
those with $S=C\cup U$ for the covers $C$.
All these inequalities define facets.
If a cover exists, the maximum positive violation is
$(q+3)/(4qN)$ for pure hypermetric and odd-clique inequalities at
$\varepsilon\widetilde d$, and $(q+2)/(4qN)$ for
\eqref{eq:cut-clique} at $d''$. These maxima are zero if there is no cover.
\end{theorem}

\begin{proof}
The block matrix $\widetilde M$ is positive definite. For integral
$z\in\Z^W$ and $w\in\Z^m$,
\begin{equation}\label{eq:orthogonal-lift}
 g_{\widetilde M}(z,w)=g_M(z)+\tau^2\sum_{j=1}^m w_j(w_j-1).
\end{equation}
The added term is nonnegative. By \cref{thm:exact-cover}, a hypermetric
violation is therefore possible exactly when
$z=(\one_C,-1)$ for a cover $C$ and
$\sum_jw_j(w_j-1)<4q$.
Its coefficient vector is
\begin{equation}\label{eq:lifted-hypermetric}
 b=(2-q-\one\trans w,\one_C,-1,w).
\end{equation}
Its gonality is at least $2q-1$, because
\[
 \abs{2-q-\one\trans w}+\sum_j\abs{w_j}\geq q-2.
\]

Write $s_M=\delta\trans M^{-1}\delta$.
The extended Schur complement parameter is
$s_{\widetilde M}=s_M+m\tau^2$. By \cref{lem:scaling},
\[
 \varepsilon s_{\widetilde M}<\frac12+\frac1{8N}<\frac34.
\]
Consequently both $\widetilde Y_\varepsilon$ and
$\widetilde Y_\varepsilon-e_0e_0\trans/4$ are positive definite.
By \eqref{eq:congruence},
$Z(\varepsilon\widetilde d)=L\widetilde Y_\varepsilon L\trans$.
If $b=(b_0,z)$, then $L\trans b=(\sig(b),-2z)$.
Hence $b\trans Zb>\sig(b)^2/4$ for nonzero $b$.
This excludes rounded psd violations with $\abs{\sig(b)}\geq3$.
The remaining rounded psd violators have coefficient sum $\pm1$ and
are the vectors \eqref{eq:lifted-hypermetric} and their negatives.
Thus none has gonality at most $2q-3$.
For $q\geq3$, triangle inequalities have gonality three and perimeter
inequalities have coefficient sum three, so the extended point is metric
feasible. Its $Z$ matrix is positive definite, and its unit diagonal
implies that all distances lie strictly between zero and one.

For a pure hypermetric violator, let $a$ entries of $w$ be $+1$ and $k$
be $-1$. Its root coefficient has absolute value at most one, so
$k-a\geq q-3=m-1$, while $k+a\leq m$. Therefore $a=0$ and
$k\in\{m-1,m\}$. This gives precisely \eqref{eq:pure-violators}.
All these vectors meet the bound $\sum_jw_j(w_j-1)<4q$.
The violations are
\[
 Q(b^*,\varepsilon\widetilde d)=\frac{q+2}{4qN},\qquad
 Q(b^{(j)},\varepsilon\widetilde d)=\frac{q+3}{4qN}.
\]
Every odd-clique violator must have coefficient sum $\pm1$, by the
preceding exclusion, so its global sign can be chosen to make it pure.
This proves the exact lists and their maximum.

Let $\Sigma$ be diagonal with entries $-1$ on $U$ and $+1$ elsewhere.
Then $Z(d'')=\Sigma Z(\varepsilon\widetilde d)\Sigma$.
The map $b\mapsto\Sigma b$ preserves odd parity, gonality, the
$\{0,\pm1\}$ restriction, and rounded psd violation.
For the vector \eqref{eq:lifted-hypermetric}, its switched vector is
\[
 \Sigma b=(2-q-\one\trans w,\one_C,1,-w).
\]
It has all nonzero entries positive exactly when
$w\in\{0,-1\}^m$ and $2-q-\one\trans w\in\{0,1\}$.
If $k$ entries of $w$ are $-1$, these conditions give
$2-q+k\in\{0,1\}$ and $k\leq q-2$, so $k=q-2$.
Thus $w=-\one_m$, the root entry is zero, and the support is $C\cup U$.
The globally negated vectors cannot give an all-positive support:
their coordinates on the nonempty cover $C$ remain negative after
switching on $U$.
Conversely, this choice is violated. Its support size is $2q-1$ and its
violation is $(q+2)/(4qN)$, establishing the exact all-positive list.

Every coefficient vector just listed has odd support size $2q-1\geq5$.
An all-positive clique inequality on an odd support of size $t$ not
containing the root is the covariance image of
the clique specialization of \eqref{eq:finite-BH} with $s=(t-1)/2$.
If its support contains the root, it is the covariance image with
$t-1$ variables and $s=(t-1)/2$.
Both meet the range in \cref{lem:clique-facets} when $t\geq5$.
The covariance map is a linear bijection of the Boolean quadric and
cut polytopes, as is seen directly on their binary vertices and convex
hulls. It preserves facets, as do cut switchings.
This proves all claimed facet properties, including those of the
maximizing vectors $b^{(j)}$.

The entries of $8qN\,\widetilde Y_\varepsilon$ are integral and lie
between zero and $8qN$. The distances and their switchings likewise
have denominator dividing $8qN$ and lie in $(0,1)$, so their numerators
are bounded by $8qN$. Switching preserves the metric polytope,
elliptope interior, and absence of rounded psd violations of gonality
at most $2q-3$; its moment matrix is positive definite by the inverse
congruence with $L$. Its principal covariance Gram matrix is therefore
positive definite as well.
The construction has polynomial size. Taking $q\geq K$ proves strong
hardness with each stated promise. Membership in $\NP$ is immediate:
the certificates have entries in $\{0,\pm1\}$ or specify a subset.
\end{proof}

\begin{remark}
If the all-positive clique family includes even supports with right-hand
side $\lfloor\abs S^2/4\rfloor$, those members are psd inequalities and
hold at $d''$. The hardness result therefore holds under either convention.
At $\varepsilon\widetilde d$, the maximum positive hypermetric and
rounded psd violation, conditional on a cover existing, is $1/(2N)$:
in \eqref{eq:orthogonal-lift}, the added term vanishes exactly when
$w\in\{0,1\}^m$. These unnormalized maxima describe the construction;
they do not give a lower bound on the violation needed for hardness.
\end{remark}
